\begin{figure}[ht]
\centering
\begin{tikzpicture}[x=1.0cm, y=1.0cm, font=\footnotesize]

% Left: complex z-plane with the unit circle and a pole pair close to it,
% showing the exact pole and the quantised pole. Right: the same pair seen
% through the direct-form coefficient grid that forces the shift.

% ----- Left panel: pole location near the unit circle -----
\begin{scope}[shift={(0,0)}]
  \draw[->, thick] (-2.3,0) -- (2.3,0) node[right] {$\Re z$};
  \draw[->, thick] (0,-2.3) -- (0,2.3) node[above] {$\Im z$};
  % unit circle
  \draw[enggray, thick] (0,0) circle (1.8);
  \node[enggray, anchor=south west, font=\scriptsize] at (1.27,1.27)
    {$\abs{z}=1$};
  % exact pole just inside (radius 0.97 -> 1.746 in figure units), angle 40 deg
  \fill[engblue] ({1.746*cos(40)},{1.746*sin(40)}) circle (2.2pt);
  \node[engblue, anchor=south, font=\scriptsize]
    at ({1.746*cos(40)},{1.746*sin(40)+0.12}) {exact pole};
  % quantised pole pushed outward past the circle (radius 1.01 -> 1.818)
  \fill[engred] ({1.818*cos(40)},{1.818*sin(40)}) circle (2.2pt);
  \node[engred, anchor=west, font=\scriptsize]
    at ({1.818*cos(40)+0.12},{1.818*sin(40)}) {quantised};
  \node[font=\scriptsize] at (0,-2.7) {(a) pole pair near $\abs{z}=1$};
\end{scope}

% ----- Right panel: impulse response, stable vs unstable -----
\begin{scope}[shift={(6.4,-1.6)}]
  \draw[->, thick] (-0.2,0) -- (4.4,0) node[right] {$n$};
  \draw[->, thick] (0,-0.2) -- (0,3.6) node[above] {$\abs{h_{n}}$};
  % stable decaying envelope (exact pole inside)
  \draw[engblue, thick] plot[smooth, domain=0:4, samples=40]
    ({\x},{2.6*exp(-0.7*\x)+0.05});
  \node[engblue, anchor=west, font=\scriptsize] at (2.2,0.75)
    {stable (inside)};
  % unstable growing envelope (quantised pole outside)
  \draw[engred, thick] plot[smooth, domain=0:3.0, samples=40]
    ({\x},{0.4*exp(0.65*\x)});
  \node[engred, anchor=south, font=\scriptsize] at (2.4,2.9)
    {unstable (outside)};
  \node[font=\scriptsize] at (2.1,-0.9) {(b) impulse response};
\end{scope}

\end{tikzpicture}
\caption[Coefficient quantisation moves an IIR filter pole across the unit
circle]{(a) A recursive digital filter with a pole pair placed just inside the
unit circle is stable; quantising the filter coefficients to a finite word
length moves the realised pole, and when the pole is close to the circle a
small coefficient perturbation can push it outside. (b) A pole inside the circle
gives a decaying impulse response; a pole outside gives one that grows without
bound, so a filter correct in exact arithmetic becomes unstable once its
coefficients are rounded. The sensitivity of the pole location to the
coefficient perturbation is the conditioning of the coefficient-to-root map of
section 16.2, and it is the reason a sharp filter is realised as a cascade of
second-order sections rather than one high-order direct form. Source: authors'
schematic.}
\label{fig:vol02ch16:iir-quantization}
\end{figure}
